\documentclass[12pt]{amsart}
\usepackage{graphicx,color,url,setspace}
\usepackage{fullpage}
\definecolor{csugreen}{cmyk}{0.92,0.18,0.94,0.61}

\usepackage{helvet}
% the following are not needed, but select particular text fonts
%\usepackage{MnSymbol}
%\usepackage[mathlf,textlf,minionint]{MinionPro}
% end optional fonts
\pagestyle{empty}
\parindent=0cm
\begin{document}
\noindent\includegraphics[width=7.5cm]{csuwrdmrk357.pdf}\\
{\begin{spacing}{0.95}\sf\footnotesize
 Colin Roberts\\
 Ph.D. Candidate\\
{\color{csugreen}DEPARTMENT OF MATHEMATICS}\\
Louis R. Weber Building\\
Campus Delivery 1874\\
Fort Collins, CO 80523-1874\\
\url{colinroberts.net}\\
\url{robertsp@rams.colostate.edu}\\
\end{spacing}
}
\smallskip
{\color{csugreen}\rule{12mm}{0.8mm}}\medskip

\today
\bigskip


\noindent Dear Editors,
\bigskip

\parindent=.5cm

I would like to submit the paper ``A Gelfand Transform for Spinor Fields on Embedded Riemannian Manifolds" to the \emph{SIAM Journal on Mathematical Analysis}. I have attached the document to this form. 

Belishev published a solution to the 2 dimensional Calder\'on in this journal in 2003. His strategy was to determine the algebra of holomorphic functions from the Dirichlet-to-Neumann operator, and then show the algebra determines the underlying manifold up to conformal class. My paper shows that in any dimension there exists a space of functions that generalizes the space of holomorphic functions and determines a compact region of $\mathbb{R}^n$ up to homeomorphism. 

To achieve this new result, we consider functions on manifolds that take values in Clifford algebras and construct a spectrum associated to this space. Along the way we build up from the basics of Clifford algebras and explain how they generalize both the complex and quaternion algebras and show that differential forms embed into Clifford valued functions. This is what allows us to bootstrap from Gelfand's result to a new class of functions.

We focus on the space of monogenic spinor fields which is well known in Clifford analysis. In dimension 2, this space is isomorphic to the algebra of complex holomorphic functions and in dimension 3 it is isomorphic to the space of quaternion harmonic fields defined by Belishev and Vakulenko in their related article which constructs a Gelfand transform for convex regions in $\mathbb{R}^3$. Hence, we also generalize their result to a much larger class of manifolds. Throughout the paper I connect the theory back to complex analysis since much of the intuition carries over nicely. 


We thank you for your consideration, and look forward to your response.



\bigskip\bigskip

\qquad\qquad Sincerely,

\qquad\qquad\qquad\qquad \includegraphics[scale=0.2]{sig.png}

\vspace{-.1in}

\qquad\qquad\qquad\qquad\qquad Colin Roberts



% \end{minipage}
\end{document}
